\chapter{La llegenda del Targus}

\epigraf{En la naturalesa no hi ha res ridícul.}{Blaise Pascal}

\lettrine[loversize=0.2, lines=2]{E}{n} les profunditats ombrívoles de Tartària, terra llegendària i remota que es perd entre les boires difuses del mite i la realitat, existia un regne antic i venerable que pocs historiadors han gosat descriure amb precisió. Era Tartària un país d’altes muntanyes coronades per neus eternes, boscos tan profunds que la llum del sol hi penetrava amb dificultat, i valls plenes d’una vegetació exuberant i estranya, on florien flors impossibles i habitaven criatures que cap home havia vist abans.

Era en aquest regne ancestral on habitava una bèstia d’aspecte formidable i grotesc anomenada pels antics savis com a Targus. Aquest monstre tenia l’aspecte d’un cuc gegantí, d’una mida tan extraordinària que el seu cos s’estenia com una enorme serp recaragolada sobre ella mateixa, més alta que els turons més prominents i més llarga que els rius més extensos. Tenia aquest cuc monstruós una pell d’una gruixudesa comparable a les muralles més impenetrables, aspra i resistent, d’una textura que ni la més afilada espasa podia penetrar amb facilitat.

Els ulls del Targus, foscos com la nit més negra i profunds com pous sense fons, irradiaven una intel·ligència maligna que infonia un terror absolut als habitants del país. Tan poderós i temible era aquest monstre que fins i tot les criatures més valentes dels mars i dels boscos de Tartària s’amagaven només sentir-ne la presència propera.

Tanmateix, la característica més sorprenent i esgarrifosa del Targus no era pas la seva aparença imponent, sinó els seus hàbits alimentaris, que horroritzaven profundament tothom que en sentia parlar. El Targus s’alimentava exclusivament del semen humà, consumint-lo en quantitats ingents i monstruoses. Aquest fet, grotesc i esfereïdor, havia estat motiu d’infinites discussions entre els savis i erudits del regne. Alguns atribuïen aquesta estranya dieta a un càstig imposat pels déus per alguna falta ancestral comesa pels habitants de Tartària; d’altres, més pragmàtics, pensaven que era simplement una adaptació biològica única, un exemple extrem de les bizarres i cruels formes que pot adoptar la natura.

Sigui com fos, la presència d’aquesta bèstia era un perill constant, una amenaça latent que mantenia el regne de Tartària en un estat permanent d’alerta i temor. Però cap història havia trasbalsat tant les ànimes dels habitants com la tragèdia que es desencadenà al castell del rei Clàmidos, governant noble i respectat d’aquest país mític. La seva filla, la princesa Seminara, famosa en tot el regne per la seva bellesa incomparable i la seva virtut immaculada, va ser víctima d’un desastre màgic sense precedents.

Un bruixot, anomenat Morvius, famós per la seva incompetència i per la seva inclinació natural cap a l’error, havia intentat realitzar un encanteri per garantir la fertilitat de les terres reials. Però, en comptes d’aconseguir-ho, va provocar un desastre: una pluja inacabable de semen màgic va caure sobre el castell, cobrint la princesa i tot el seu entorn en una substància viscosa i repugnant.

Aquest fet vergonyós i absurd arribà ràpidament a les orelles del Targus, despertant la seva fam insaciable. Amb la ferocitat d’un drac antic i la determinació d’un exèrcit invasor, la bèstia emergí de les profunditats fosques de la seva cova subterrània i es dirigí directament cap al castell.

El terror s’apoderà ràpidament de la cort reial, i ningú semblava tenir prou valentia per plantar cara al monstre que avançava inexorable. En aquest moment desesperat, aparegué el noble cavaller Sant Johan, famós per la seva força sobrehumana, la seva pietat exemplar i una valentia llegendària. Armat amb una llança forjada amb ferro extret d’un meteorit caigut del cel i beneïda per tots els savis de Tartària, aquest valent cavaller jurà defensar la princesa i l’honor del regne.

El combat que s’inicià fou llarg, intens i terriblement viscós. Cavaller i cuc lluitaren durant tres dies i tres nits sense descans. Cada embat del Targus feia tremolar la terra, cada cop de la llança del cavaller intentava en va travessar aquella pell invulnerable. El camp de batalla quedà completament cobert d’una substància repulsiva, dificultant encara més els moviments i les maniobres.

El tercer dia, quan el sol començava a desaparèixer darrere les muntanyes, semblava que la victòria era impossible i que Tartària estava condemnada. Però fou aleshores quan Sant Johan, en un últim esforç desesperat, pronuncià el nom secret i sagrat de Tartària, canalitzant tota l’energia dels ancestres que havien defensat aquella terra en temps antics, i amb una força sobrenatural, clavà la seva llança directament al cor de la bèstia.

De la ferida mortal que infligí el cavaller sorgí una cascada de sang prodigiosa que cobrí tota la vall circumdant. I, d’aquesta sang vessada, emergiren miraculosament centenars de motxilles meravelloses i extraordinàries, d’una bellesa i funcionalitat mai vistes.

Sant Johan trià la motxilla més magnífica de totes i la lliurà amb humilitat i devoció a la princesa Seminara com a símbol etern del seu rescat. Aquesta motxilla prodigiosa fou guardada amb cura extrema pel rei Clàmidos i els seus descendents, fins que, al cap de molts segles, arribà misteriosament fins a les mans de David Bagan, descendent de la princesa, qui sense ser plenament conscient del significat real d’aquesta relíquia sagrada, en fou el guardià involuntari.

I així, al llarg dels segles, la història del Targus ha estat transmesa de generació en generació, perpetuant-se com a patrimoni mític de Tartària, i explicada als seminaris i passadissos de la venerable Universitat Autònoma, on estudiants perplexos encara avui escolten bocabadats les gestes fabuloses del cosmos mitològic de la UAB.